\documentclass[10pt]{article}
\usepackage[T1]{fontenc}
\usepackage{lmodern}
\usepackage[margin=.98in,top=.76in,bottom=.78in]{geometry}
\usepackage{amsmath,amssymb,amsthm,booktabs,microtype}
\usepackage[hidelinks]{hyperref}
\newtheorem{theorem}{Theorem}
\newtheorem{lemma}[theorem]{Lemma}
\newcommand{\oeisentry}[1]{\href{https://oeis.org/#1}{\texttt{#1}}}
\newcommand{\eps}{\varepsilon}
\title{\vspace{-1.2em}Lengths and letter positions for two binary pair-replacement rules}
\author{Blake Stone mathematics repository\\[-2pt]\small Proof developed with OpenAI Codex}
\date{November 26, 2026}
\begin{document}
\maketitle
\begin{abstract}
Two nonoverlapping replacement rules proposed by Clark Kimberling have three-letter
descriptions in terms of the gaps between ones. We derive both exact iterate-length
recurrences, four algebraic position limits, and all-index bounds. The results prove
the proposed order-seven length recurrence, correct a Tribonacci sequence reference,
and isolate a single failing index in a proposed inequality.
\end{abstract}

\section{Statements}
Begin at $w_0=00$. Scan the current word from the left, replacing $00$ by $001$, replacing
each $1$ by its prescribed image, and copying any unpaired $0$. The inserted letters are first
scanned on the next iteration. Write $w_{n+1}=R(w_n)$, $L_n=|w_n|$, and use one-based
positions $p(n)$ and $q(n)$ of the $n$-th one and zero in the limiting word.
The two choices $R_T(1)=010$ and $R_C(1)=000$ give \oeisentry{A285666} and \oeisentry{A286939},
respectively. Kimberling's June 2017 entries explicitly ask about lengths and positions.
Earlier numerical matches to alternative morphisms for \oeisentry{A286058} appear in Sequence
Machine \cite{sm}. The proofs here start from the left-to-right replacement itself.

\begin{theorem}\label{th:length}
For $R_T$, $L_0=2,L_1=3,L_2=6$ and $L_{n+3}=L_{n+2}+L_{n+1}+L_n$ for $n\ge0$.
Thus $L_n=\text{\oeisentry{A001590}}(n+4)$. The proposed \oeisentry{A001580}, whose $n$-th term is
$2^n+n^2$, already gives $8$ instead of $|001010|=6$ at $n=2$.

For $R_C$ the lengths start $2,3,6,10,17,29,51$. For every $n\ge0$,
\[
 L_{n+7}=L_{n+6}+2L_{n+5}-L_{n+4}-L_{n+3}+2L_{n+2}-2L_n.
\]
This proves the identification in \oeisentry{A285665}, including its eighth term $90$.
\end{theorem}

\begin{theorem}\label{th:dens}
Let $\alpha>1$ satisfy $\alpha^3=\alpha^2+\alpha+1$, and $\beta>1$ satisfy
$\beta^3=2\beta+2$. For both rules $m_1=\lim p(n)/n$ and $m_0=\lim q(n)/n=m_1/(m_1-1)$
exist. Their values are
\[
 m_1(T)=\frac{2+3\alpha+4\alpha^2}{1+\alpha+\alpha^2},\qquad
 m_1(C)=\frac{3\beta^2+6\beta+4}{\beta^2+\beta}.
\]
In the order $(m_1,m_0)$ these are approximately $(3.382976,1.419643)$ and
$(4.899688,1.256431)$, respectively. The proposed strict unit bounds on
$q(n)/n-m_0$ hold for all $n\ge1$ in \oeisentry{A285667} and \oeisentry{A285432}. The same strict
unit bounds hold for $p(n)/n-m_1$ in \oeisentry{A286058}. For \oeisentry{A285564} the proposed
$-1<p(n)/n-m_1<2$ holds exactly for $n\ge2$. Its lower bound fails at $n=1$, where
$p(1)=3$.
\end{theorem}

\section{The substitutions on gaps}
Write $D(k)=10^k$ and extend $D$ to a word of gap symbols by concatenation. The gap
alphabets below use the letters $a,b,c$ in the order displayed.
\begin{center}
\begin{tabular}{lccccc}
\toprule
Rule & $(a,b,c)$ in zeros & $h$ & $S(a)$ & $S(b)$ & $S(c)$\\
\midrule
$R_T$ & $(1,2,3)$ & $a=1$ & $c$ & $ca$ & $cb$\\
$R_C$ & $(2,5,6)$ & $b=5$ & $b$ & $ac$ & $aab$\\
\bottomrule
\end{tabular}
\end{center}
Here $h$ denotes both the indicated letter and its zero count. For each of the six cases
one directly checks
\[
                         R(D(k))\,0^h=0^h D(S(k)).
\tag{1}
\]
The rule can be concatenated at a boundary immediately before a $1$. For example under
$R_C$, $D(5)$ rewrites to $000(001)^2 0$. Appending five zeros gives $0^5D(2)D(6)$.

The infinite gap word is explicitly
\[
                         K=h\,S(h)\,S^2(h)\,S^3(h)\cdots .
\tag{2}
\]
Indeed $K=hS(K)$. Taking longer prefixes in (1) gives $R(D(K))=0^hD(S(K))$, so
$W=00D(K)$ is fixed by $R$. The scan preserves the prefix relation: the only cut that can
split a token leaves one trailing zero, and the image of $00$ begins with zero. The iterates
of $00$ are consequently nested prefixes of $W$. Their lengths tend to infinity, because
each iterate begins with $00$ and replacing that pair increases length. This proves that
$W$ is precisely the word defined in the entries.

\section{Finite right endpoints and lengths}
Equation (1) is an identity with an appended suffix. To compute the lengths of the
unextended finite words, keep a terminal binary word $s_j$. The following table gives a
cycle of such words and a possibly empty gap string $J_j$. In each part the next row
wraps back to the first. The notation $(k_1,k_2;d)$ means $D(k_1)D(k_2)10^d$,
and $(\,;d)$ means $10^d$. All rows assert the finite identity
$R(s_j)=0^h D(J_j)s_{j+1}$, which follows by applying the original scan.
\begin{center}
\begin{tabular}{ccccc}
\toprule
$R_T$ suffix & $J_j$ & \hspace{.35in} & $R_C$ suffix & $J_j$\\
\midrule
$(\,;0)$ & $\eps$ && $(5,2;5)$ & $ac$\\
$(\,;1)$ & $\eps$ && $(5,2;1)$ & $\eps$\\
$(\,;2)$ & $c$    && $(2,6;4)$ & $baa$\\
&&& $(5,2;0)$ & $\eps$\\
&&& $(2,6;3)$ & $ba$\\
&&& $(2,5;1)$ & $\eps$\\
\bottomrule
\end{tabular}
\end{center}
At $n=1$ in the first rule and $n=4$ in the second, $w_n=00s_0$.
If $w_n=00D(v)s_j$, applying (1) and the table gives exactly
\[
                        w_{n+1}=00D\big(hS(v)J_j\big)s_{j+1}.
\tag{3}
\]
The table is also a small reproducible calculation of both recurrences. To give the
details, let $X$ be the column vector of letter counts in $v$, $e(w)$ the count column
of a word, and $d=(1+a,1+b,1+c)$ the row of tile lengths. Then
\[
 X'=MX+e(h)+e(J_j),\qquad L_n=2+dX+|s_j|,\qquad
 M_T=\begin{pmatrix}0&1&0\\0&0&1\\1&1&1\end{pmatrix},\quad
 M_C=\begin{pmatrix}0&1&2\\1&0&1\\0&1&0\end{pmatrix}.
\]
The characteristic polynomials are $x^3-x^2-x-1$ and $x^3-2x-2$. Hence each associated
third-order length residual depends only on the suffix phase: its coefficient of $X$
vanishes by Cayley--Hamilton. For $R_T$, evaluating the residual
$L_{n+3}-L_{n+2}-L_{n+1}-L_n$ at $n=1,2,3$ gives $(0,0,0)$. The initial lengths check
$n=0$. For $R_C$, the residual $D_n=L_{n+3}-2L_{n+1}-2L_n$ repeats in six phases from
$n=4$. With direct evaluation also at $0\le n<4$, its cycle from $n=0$ is
\[
                          (0,-1,-3,-3,-2,0).
\]
These six numbers satisfy $D_{n+4}-D_{n+3}+D_{n+1}-D_n=0$, including at the wrap.
Multiplying $x^3-2x-2$ by $x^4-x^3+x-1$ gives the recurrence in
Theorem~\ref{th:length}. This proves the length statements.

\section{Frequencies and bounds}
Both matrices are primitive: their directed graphs are strongly connected, with a loop
in the first and cycles of lengths two and three in the second. Perron--Frobenius gives
limiting normalized count vectors for every $S^j(k)$, uniformly over the three choices of
$k$. This conclusion applies to all prefixes of $K$: for fixed $j$ discard
$h\cdots S^{j-1}(h)$, and decompose its infinite tail into complete $S^j$-images.
At most one end tile is cut, with length bounded for fixed $j$. First take the prefix
length to infinity and then $j$ to infinity. The right eigenvectors, in our order of
letters, are proportional to $(1,\alpha,\alpha^2)$ and $(\beta^2-1,\beta,1)$.
Since $p(1)=3$ and the successive gaps are $1+k$ for $k$ in $K$, the average tile
length gives the two displayed $m_1$ values. Inverting the increasing one-position
function gives density $1/m_1$ of ones in binary prefixes. The complementary zero
density is $1-1/m_1$, proving the formula for $m_0$. For use below, exact substitution
in the cubics brackets $1.8392<\alpha<1.8393$ and $1.7692<\beta<1.7693$.
Inserting these rational intervals in the displayed formulae gives
$3<m_1(T)<7/2$ and $4<m_1(C)<5$.

Both binary words have isolated ones and start with a zero, hence $1\le q(n)/n\le2$
for all $n$. Both $m_0$ lie strictly between 1 and 2, which proves the two zero bounds.

For $R_T$, every image begins in $c$ and each $a$ in an image is terminal.
From $K=aS(K)$, two $a$'s are never adjacent. The position $p(n)$ is a sum of $n$
increments, the first equal to $3$ and the rest between $2$ and $4$, with no adjacent
$2$'s. At most $\lfloor n/2\rfloor$ increments equal $2$. Thus
$5/2\le p(n)/n\le4$. Together with $3<m_1(T)<7/2$ this gives the strict inequalities.

For $R_C$ we prove $4\le p(n)/n<6$ for $n\ge2$. For the lower bound, subtract $4$ from
each increment. The first contributes $-1$. The successive gap letters contribute
$-1,2,3$ for $a,b,c$ respectively. From $K=bS(K)$ its first three letters are $bac$.
There are never more than two successive $a$'s since every image ends in $b$ or $c$.
The first one or two letters of $K$ have sum at least one. The suffix starting with the
third letter can be grouped into a $b$ or $c$ followed by at most two $a$'s. Every
prefix of each such group has nonnegative sum. This proves the lower bound. For the
upper bound, subtract $6$ from each increment. Letter weights are $-3,0,1$, and each
$c$ is immediately preceded by $a$, since it appears only in the image $ac$. Pairing
these letters bounds every prefix sum by zero, in addition to the initial $3-6$.
This proves the claimed upper bound even at $n=1$. Since $4<m_1(C)<5$,
the proposed lower bound fails there and holds at every later index.
\hfill$\square$

\paragraph{Verification.}
The accompanying Lean file proves both recurrences on the literal finite binary words,
including the boundary identities and initial cases. The analytic limits and inequalities
are proved in this manuscript. An independent verifier implements the original scan,
checks the suffix tables and recurrence against literal words, and checks the position
bounds and the failed first index.

\begin{thebibliography}{9}
\bibitem{oeis} C. Kimberling, entries \oeisentry{A285665}, \oeisentry{A285666}, \oeisentry{A285667},
\oeisentry{A286058}, \oeisentry{A286939}, \oeisentry{A285432} and \oeisentry{A285564}, June 2017,
\emph{The On-Line Encyclopedia of Integer Sequences}. The comparison sequences are
\oeisentry{A001580} and \oeisentry{A001590}.
\bibitem{sm} \emph{Sequence Machine}, numerical matches for A286058.\newline
\url{https://sequencedb.net/index.html?p=600&s=diff}.
\end{thebibliography}
\end{document}
